
\documentclass[addpoints]{exam}
\usepackage{amsmath}
\usepackage{amssymb}
\usepackage{color}

% \printanswers

\newcommand{\event}{Final Exam}
\newcommand{\tournament}{}
\def\type{0}

\setlength\fillinlinelength{\textwidth}
\CorrectChoiceEmphasis{\color{red}\bfseries}
\SolutionEmphasis{\color{red}\bfseries}

\setlength\answerclearance{2pt}
\pagestyle{headandfoot}
\runningheadrule
\runningheader{\event}{\tournament}{\ifnum\type=1 Class Set Test \else Full Name:\hspace{1cm} \fi}
\runningfootrule
\runningfooter{}{Page \thepage\ of \numpages}{}

\begin{document}
\begin{center}
    {\huge Grade 10 Principles of Mathematics \\ \event
    \ifprintanswers \space Key
    \else \space \ifnum\type<2 Test \fi \ifnum\type=1 \textbf{(CLASS SET)} \fi
          \ifnum\type=2 Answer Sheet \fi
    \fi \\ \vspace{3mm}\tournament\vspace{1mm}}
\end{center}

\vspace{.25\textheight}

\begin{center}
    First Name: \ifprintanswers
    \underline{\hspace{3.5cm}}\textbf{KEY}\underline{\hspace{5.5cm}}\\\vspace{5mm}
    \else \underline{\hspace{10cm}}\\ \vspace{5mm} \fi
    Last Name: \underline{\hspace{10cm}}\\ \vspace{5mm}
\end{center}

\noindent\textbf{\underline{Directions:}}
\begin{itemize}
    \ifnum\type=1 \item \textbf{This test is a class set.} Do not write on this paper. \fi
    \item Show all work for full marks. Leave answers in exact form where appropriate.
    \item The exam is designed to be completed in 75 minutes.
\end{itemize}
\begin{center}
\textit{For grading use only}\\
\multirowgradetable{1}[pages]
\end{center}

\newpage
\begin{questions}

\section*{Part A --- Multiple Choice (20 marks)}
\vspace{5mm}

%--- Linear Systems (4) ---
\question[1] The solution to the system $y = 2x + 1$ and $y = -x + 7$ is:
\begin{choices}
    \choice $(2, 5)$
    \correctchoice $(2, 5)$
    \choice $(3, 4)$
    \choice $(1, 6)$
\end{choices}

\question[1] Which system has \textbf{no solution}?
\begin{choices}
    \correctchoice $y = 3x + 1$ and $y = 3x - 4$
    \choice $y = 2x + 1$ and $y = -2x + 1$
    \choice $x + y = 5$ and $2x + 2y = 10$
    \choice $y = x$ and $y = -x$
\end{choices}

\question[1] A system is solved by elimination. Adding the equations $2x + 3y = 11$ and
$-2x + y = 1$ gives:
\begin{choices}
    \choice $4y = 12$
    \correctchoice $4y = 12$
    \choice $2y = 10$
    \choice $4x = 10$
\end{choices}

\question[1] Two numbers have a sum of 40 and a difference of 8. The larger number is:
\begin{choices}
    \choice $16$
    \choice $20$
    \correctchoice $24$
    \choice $32$
\end{choices}

%--- Analytic Geometry (4) ---
\question[1] The midpoint of the segment joining $(-2, 4)$ and $(6, -2)$ is:
\begin{choices}
    \choice $(4, 1)$
    \correctchoice $(2, 1)$
    \choice $(2, 3)$
    \choice $(4, 2)$
\end{choices}

\question[1] The distance between $(1, -3)$ and $(5, 0)$ is:
\begin{choices}
    \choice $3$
    \choice $4$
    \correctchoice $5$
    \choice $7$
\end{choices}

\question[1] The slope of a line perpendicular to $y = \dfrac{2}{3}x - 4$ is:
\begin{choices}
    \choice $\dfrac{2}{3}$
    \choice $-\dfrac{2}{3}$
    \correctchoice $-\dfrac{3}{2}$
    \choice $\dfrac{3}{2}$
\end{choices}

\question[1] Which point lies on the circle $x^2 + y^2 = 25$?
\begin{choices}
    \choice $(2, 4)$
    \correctchoice $(3, 4)$
    \choice $(4, 4)$
    \choice $(5, 1)$
\end{choices}

%--- Quadratic Expressions (4) ---
\question[1] Expand and simplify: $(x + 3)(x - 5)$
\begin{choices}
    \choice $x^2 - 2x - 15$
    \correctchoice $x^2 - 2x - 15$
    \choice $x^2 + 2x - 15$
    \choice $x^2 - 8x - 15$
\end{choices}

\question[1] Factor completely: $2x^2 - 8$
\begin{choices}
    \choice $2(x^2 - 8)$
    \choice $(2x - 4)(x + 2)$
    \correctchoice $2(x-2)(x+2)$
    \choice $2(x-4)(x+1)$
\end{choices}

\question[1] Factor: $6x^2 + 7x - 3$
\begin{choices}
    \choice $(6x - 1)(x + 3)$
    \correctchoice $(3x - 1)(2x + 3)$
    \choice $(2x - 1)(3x + 3)$
    \choice $(3x + 1)(2x - 3)$
\end{choices}

\question[1] The expression $x^2 - 10x + 25$ factors as:
\begin{choices}
    \choice $(x-5)(x+5)$
    \choice $(x-5)(x-4)$
    \correctchoice $(x-5)^2$
    \choice $(x+5)^2$
\end{choices}

%--- Quadratic Functions (4) ---
\question[1] The vertex of $y = 2(x - 3)^2 + 1$ is:
\begin{choices}
    \choice $(-3, 1)$
    \correctchoice $(3, 1)$
    \choice $(3, -1)$
    \choice $(-3, -1)$
\end{choices}

\question[1] The parabola $y = -x^2 + 4x - 3$ opens:
\begin{choices}
    \correctchoice Downward, with a maximum value
    \choice Upward, with a minimum value
    \choice Downward, with a minimum value
    \choice Upward, with a maximum value
\end{choices}

\question[1] A ball follows $h(t) = -5t^2 + 20t + 1$. Its maximum height occurs at $t =$:
\begin{choices}
    \choice $1$ s
    \choice $3$ s
    \correctchoice $2$ s
    \choice $4$ s
\end{choices}

\question[1] The discriminant of $2x^2 - 3x + 5 = 0$ is:
\begin{choices}
    \choice $49$
    \choice $-31$
    \correctchoice $-31$
    \choice $31$
\end{choices}

%--- Quadratic Equations & Trig (4) ---
\question[1] Solve $x^2 - 5x + 6 = 0$. The solutions are:
\begin{choices}
    \choice $x = 1$ or $x = 6$
    \choice $x = -2$ or $x = -3$
    \correctchoice $x = 2$ or $x = 3$
    \choice $x = -1$ or $x = 6$
\end{choices}

\question[1] In a right triangle with opposite $= 7$ and hypotenuse $= 25$, $\sin\theta =$:
\begin{choices}
    \choice $\dfrac{25}{7}$
    \correctchoice $\dfrac{7}{25}$
    \choice $\dfrac{24}{25}$
    \choice $\dfrac{7}{24}$
\end{choices}

\question[1] A ladder 10 m long leans against a wall at an angle of $65^\circ$ to the ground.
The height it reaches is approximately:
\begin{choices}
    \choice $4.2$ m
    \choice $7.7$ m
    \correctchoice $9.1$ m
    \choice $10.6$ m
\end{choices}

\question[1] Using the quadratic formula on $x^2 + 2x - 8 = 0$, the solutions are:
\begin{choices}
    \choice $x = 1$ or $x = -8$
    \correctchoice $x = 2$ or $x = -4$
    \choice $x = -2$ or $x = 4$
    \choice $x = 4$ or $x = -1$
\end{choices}


\section*{Part B --- Short Answer (40 marks)}

\question[8] \textbf{Linear Systems.} Solve the system using elimination and verify your solution.
Then set up and solve a system for the following problem: two numbers have a sum of 54 and the
larger is 12 more than twice the smaller. Find both numbers.
\begin{parts}
    \part[4] Solve by elimination: $3x + 2y = 16$ and $5x - 2y = 8$.
    \part[4] Word problem: set up and solve.
\end{parts}
\begin{solution}[10cm]
\textbf{(a)} Adding: $8x = 24 \Rightarrow x = 3$

Substituting: $3(3) + 2y = 16 \Rightarrow 2y = 7 \Rightarrow y = 3.5$

\textbf{Solution: $(3, 3.5)$}

Verify: $5(3) - 2(3.5) = 15 - 7 = 8$ \checkmark

\textbf{(b)} Let $x$ = smaller, $y$ = larger.
\[x + y = 54 \quad \text{and} \quad y = 2x + 12\]
Substituting: $x + 2x + 12 = 54 \Rightarrow 3x = 42 \Rightarrow x = 14$, $y = 40$

\textbf{The numbers are 14 and 40.}
\end{solution}

\question[8] \textbf{Analytic Geometry.} The vertices of triangle $ABC$ are $A(0, 6)$, $B(-4, 0)$, and $C(4, 0)$.
\begin{parts}
    \part[3] Find the length of each side and classify the triangle.
    \part[3] Find the equation of the median from $A$ to the midpoint of $BC$.
    \part[2] Find the equation of the altitude from $B$ to side $AC$.
\end{parts}
\begin{solution}[10cm]
\textbf{(a)}
\[AB = \sqrt{(-4-0)^2+(0-6)^2} = \sqrt{16+36} = \sqrt{52} = 2\sqrt{13}\]
\[AC = \sqrt{(4-0)^2+(0-6)^2} = \sqrt{16+36} = 2\sqrt{13}\]
\[BC = \sqrt{(4-(-4))^2+0} = 8\]

$AB = AC$, so the triangle is \textbf{isosceles}.

\textbf{(b)} Midpoint of $BC$: $M = \left(\dfrac{-4+4}{2}, \dfrac{0+0}{2}\right) = (0, 0)$

Slope $AM = \dfrac{6-0}{0-0}$ is undefined — $AM$ is the $y$-axis.

Equation of median: $\mathbf{x = 0}$

\textbf{(c)} Slope of $AC = \dfrac{0-6}{4-0} = -\dfrac{3}{2}$

Slope of altitude from $B \perp AC$: $m = \dfrac{2}{3}$

Through $B(-4, 0)$: $y - 0 = \dfrac{2}{3}(x + 4)$
\[\mathbf{y = \dfrac{2}{3}x + \dfrac{8}{3}}\]
\end{solution}

\question[8] \textbf{Quadratic Expressions.} Factor each expression completely.
\begin{parts}
    \part[2] $4x^2 - 36y^2$
    \part[2] $3x^2 + 9x - 30$
    \part[2] $6x^2 - x - 12$
    \part[2] $2x^3 - 8x^2 + 6x$
\end{parts}
\begin{solution}[8cm]
\textbf{(a)} $4(x^2 - 9y^2) = \mathbf{4(x-3y)(x+3y)}$

\textbf{(b)} $3(x^2 + 3x - 10) = \mathbf{3(x+5)(x-2)}$

\textbf{(c)} Product $= 6 \times (-12) = -72$; need factors of $-72$ that add to $-1$: $8$ and $-9$.
\[6x^2 + 8x - 9x - 12 = 2x(3x+4) - 3(3x+4) = \mathbf{(2x-3)(3x+4)}\]

\textbf{(d)} $2x(x^2 - 4x + 3) = 2x(x-1)(x-3) = \mathbf{2x(x-1)(x-3)}$
\end{solution}

\question[8] \textbf{Quadratic Functions.} A ball is thrown upward from a height of 1 m.
Its height is given by $h(t) = -5t^2 + 30t + 1$, where $h$ is in metres and $t$ in seconds.
\begin{parts}
    \part[3] Complete the square to write $h(t)$ in vertex form.
    \part[2] State the maximum height and when it occurs.
    \part[3] Find when the ball hits the ground (to 2 decimal places).
\end{parts}
\begin{solution}[10cm]
\textbf{(a)} $h(t) = -5(t^2 - 6t) + 1 = -5(t^2 - 6t + 9 - 9) + 1$
\[= -5(t-3)^2 + 45 + 1 = \mathbf{-5(t-3)^2 + 46}\]

\textbf{(b)} Vertex $(3, 46)$: maximum height is $\mathbf{46}$ \textbf{m} at $\mathbf{t = 3}$ \textbf{s}.

\textbf{(c)} Set $h(t) = 0$: $-5t^2 + 30t + 1 = 0 \Rightarrow 5t^2 - 30t - 1 = 0$
\[t = \frac{30 \pm \sqrt{900 + 20}}{10} = \frac{30 \pm \sqrt{920}}{10}\]
$\sqrt{920} \approx 30.33$, so $t = \dfrac{30 + 30.33}{10} \approx \mathbf{6.03}$ \textbf{s}

(negative root rejected since $t \geq 0$)
\end{solution}

\question[8] \textbf{Trigonometry.} A surveyor at point $A$ observes the top of a building at an angle of
elevation of $38^\circ$. She walks 50 m closer to the base of the building to point $B$ and observes the
top at an angle of elevation of $62^\circ$. Both points $A$ and $B$ are on level ground.
\begin{parts}
    \part[4] Let $d$ be the distance from $B$ to the base of the building and $h$ be the height.
    Write two equations relating $h$ and $d$.
    \part[4] Solve for $h$ (the height of the building) to the nearest metre.
\end{parts}
\begin{solution}[10cm]
\textbf{(a)} From $B$: $\tan 62^\circ = \dfrac{h}{d} \Rightarrow h = d\tan 62^\circ$

From $A$: $\tan 38^\circ = \dfrac{h}{d + 50} \Rightarrow h = (d+50)\tan 38^\circ$

\textbf{(b)} Setting equal:
\[d\tan 62^\circ = (d+50)\tan 38^\circ\]
\[d(1.8807) = (d+50)(0.7813)\]
\[1.8807d = 0.7813d + 39.065\]
\[1.0994d = 39.065 \Rightarrow d = 35.53 \text{ m}\]

$h = 35.53 \times \tan 62^\circ = 35.53 \times 1.8807 \approx \mathbf{66.8 \text{ m}}$

\textbf{The building is approximately 67 m tall.}
\end{solution}

\end{questions}
\end{document}
